%==============================================================================
\chapter{الميزانية والموارد}
%==============================================================================

\section{الإطار}
يندرج المشروع في إطار \textbf{أكاديمي (مشروع مصغّر)}. لذلك تُترجَم الميزانية
أساسًا إلى \textbf{موارد بشرية وبرمجية وعتادية} مُعبَّأة، دون كلفة مالية مباشرة.

\section{الموارد المُعبَّأة}
\begin{xltabular}{\textwidth}{>{\bfseries}p{4cm} X}
\toprule
\textbf{المورد} & \textbf{التفصيل} \\
\midrule
\endhead
بشرية & فريق مشروع طلابي (التصميم، التطوير، الاختبار، التوثيق). \\
برمجية & بيئة تنفيذ بالحاويات، نظام إدارة قواعد معطيات علائقي، بيئة تطوير، أدوات UML، سلسلة التصريف. \\
عتادية & محطّات تطوير، مستشعرات مُحاكاة لشقّ المؤشّرات. \\
خارجية & حساب OAuth Google، شهادات X.509 (موقّعة ذاتيًا في العرض). \\
\bottomrule
\end{xltabular}

\begin{encadre}
\textbf{قيد أكاديمي كبير:} أي استعمال لمولّدات الشيفرة (ChatGPT، DeepSeek
ومشتقّاتها) ممنوع ويُعتبر غشًّا في اختبار الامتحان.
\end{encadre}

%------------------------------------------------------------------------------
\section{التقييم الاقتصادي: فرضية التصنيع}
%------------------------------------------------------------------------------
يصف القسم السابق الكلفة \emph{الفعلية} للمشروع، وهي معدومة من حيث الإنفاق
المباشر. وعلى سبيل التقييم، يُقدّر هذا القسم ما سيكلّفه المنتج نفسه \textbf{لو
أنجزته شركة خدمات برمجية}. والمبالغ المذكورة \textbf{فرضيات عمل} لا نفقات
ملتزَمًا بها.

\subsection{عبء العمل المرجعي}
يعتمد مخطّط العبء في فصل الآجال \textbf{304 نقطة قصّة} على 8 دورات. فبمعدّل
$0{,}79$ يوم-رجل للنقطة، يبلغ العبء $\approx 240$ يوم-رجل؛ وبمعدّل $1$ يوم-رجل
للنقطة --- وهي فرضية متحفّظة لفريق يكتشف المنظومة --- يبلغ $\approx 304$
يوم-رجل. وهذا المجال هو أساس التقدير.

\subsection{توزيع غلاف مرجعي}
لغلاف قدره \textbf{200\,000 دولار أمريكي}، يعكس التوزيع التالي ممارسات السوق.
وحصّة التطوير ليست الأغلبية: ومن الأخطاء الشائعة في التقدير مطابقتها بالكلفة
الإجمالية.

\begin{xltabular}{\textwidth}{>{\bfseries}p{5.2cm} r r X}
\toprule
\textbf{البند} & \textbf{\%} & \textbf{USD} & \textbf{المبرّر} \\
\midrule
\endhead
التطوير & 45 & 90\,000 & الواجهة الخلفية والأمامية والتكامل \\
إدارة المشروع والتحليل & 12 & 24\,000 & البند الأكثر بخسًا في التقدير \\
الجودة والاختبار & 10 & 20\,000 & تغطية $\geq 70\,\%$، اختبارات تكامل وحِمل \\
تصميم التجربة والنفاذية & 8 & 16\,000 & كلفة مخفّضة: الميثاق التصميمي قائم \\
البنية التحتية والأدوات (السنة 1) & 6 & 12\,000 & الاستضافة والنسخ الاحتياطي والإشراف \\
تدقيق أمني خارجي & 5 & 10\,000 & معطيات تموضع جغرافي حسّاسة \\
احتياطي للطوارئ & 14 & 28\,000 & بدون احتياطي، ينحرف غلاف بهذا الحجم \\
\bottomrule
\end{xltabular}

\subsection{الحساسية للتعريفة اليومية}
العامل المهيمن جغرافيّ لا تقنيّ. فبتعريفة متوسّطة تتراوح بين 600 و900 دولار
لليوم-رجل، يموّل الغلاف $\approx 220$--$280$ يوم-رجل: أي أنّه يغطّي العبء
المرجعي بالكاد، \textbf{دون هامش}. وبتعريفة تتراوح بين 250 و400 دولار، يموّل
$\approx 500$--$700$ يوم-رجل، ويتيح فريقًا من خمسة إلى ستّة أشخاص مع احتياطي
سليم.

\subsection{المجال القابل للتعديل}
إذا ضاق الغلاف، فإنّ ترتيب الحذف التالي يحفظ جوهر المنتج: أوّلًا الوظائف
المتقدّمة (الملحق: تتبّع الملكة، رمز QR للدفعة)، ثمّ الكشف التكيّفي عن الشذوذ
--- الذي تمثّل معايرته دون سجلّ تاريخي حقيقي خطرًا محدّدًا ---، وأخيرًا الاستقبال
الفعلي لمعطيات المستشعرات، إذ يكفي نموذج المعطيات في الإصدار الأوّل. وتبقى غير
قابلة للتقليص: الكيانات المهنية، والزيارات، والاستيثاق والتحكّم في النفاذ، والنمط
غير المتّصل، ولوحات القيادة.

\subsection{الكلفة الإجمالية للتملّك}
لا تمثّل كلفة البناء سوى جزء من كلفة حياة المنتج. فالصيانة التصحيحية والتطويرية
تُقدَّر عادةً بـ\textbf{15--20\,\% سنويًا} من غلاف البناء، أي 30\,000--40\,000
دولار في السنة. وعلى ثلاث سنوات، تقترب الكلفة الإجمالية للتملّك من
\textbf{300\,000 دولار}. ويُضاف إليها ما يُغفَل عادةً في التقدير الأوّلي: استرجاع
المعطيات القائمة لدى المستثمِر، والتكوين ومرافقة التغيير لدى مستعملي الميدان،
والترجمة الاحترافية المراجَعة للّغات الثلاث، وأعمال المطابقة المتّصلة بالمعطيات
ذات الطابع الشخصي.

\begin{encadre}
\textbf{نطاق هذا القسم:} يتعلّق الأمر بـ\emph{تقييم} بيداغوجي يهدف إلى وضع جهد
المشروع ضمن رتب المقادير السائدة في السوق. وهو لا يلزم بأيّ إنفاق ولا يغيّر
الإطار الأكاديمي المحدَّد في مستهلّ الفصل.
\end{encadre}
